\thispagestyle{empty}
\vspace*{18mm}
\begin{center}
  {\sffamily\small ENGINEERING \& DATA SCIENCES \,$\cdot$\, CHEMISTRY S1-1}\\[10pt]
  {\sffamily\bfseries\Huge\color{heading} From atomic entities}\\[4pt]
  {\sffamily\bfseries\Huge\color{heading} to physico-chemical systems}\\[12pt]
  {\sffamily\LARGE \ifkey Worksheet answers\else Worksheets\fi}\\[4pt]
  {\large 5 units \,$\cdot$\, 15 questions each, from recall to challenge}\\[3pt]
  {\normalsize Semester 1 \,$\cdot$\, Academic Year 2026--2027}
\end{center}
\vspace{10mm}
{\sffamily\bfseries\color{heading} How to use \ifkey the answers\else the worksheets\fi}
\begin{itemize}[leftmargin=7mm, itemsep=3pt]
\ifkey
  \item Each question is followed by a full worked answer in a \textcolor{keyblue}{blue box}. Compare
        your \emph{method} as well as your result: units, significant figures, justification.
  \item If a Level 1 or 2 question goes wrong, re-read the matching part of the study guide before
        moving on; the Level 3 and Challenge questions build on them.
\else
  \item Revise the unit (study guide: concepts and recap) before starting its worksheet.
  \item Work on paper, without looking at the answers. Write the formula with letters first, then the
        numbers with their units.
  \item The questions get harder as you go. Aim for all of Levels 1--2 before attempting Level 3; the
        Challenge questions link several ideas or go slightly beyond the handout.
  \item Check yourself with \emph{Chem-S1-1\_Worksheet-Answers.pdf}.
\fi
\end{itemize}
\vspace{4mm}
{\sffamily\bfseries\color{heading} Levels}\par\vspace{1mm}
\begin{tabular}{@{}l l@{}}
\level{1} & Q1--4: recall a definition, apply one formula.\\
\level{2} & Q5--8: explain a result or chain two steps.\\
\level{3} & Q9--12: multi-step problems, interpreting data.\\
\level{4} & Q13--15: synthesis, links between units, going beyond the handout.
\end{tabular}
\vspace{6mm}

{\sffamily\bfseries\color{heading} Contents}\par\vspace{2mm}
\begin{center}
\renewcommand{\arraystretch}{1.3}
\begin{tabular}{@{}c l@{}}
\toprule
Worksheet & Unit\\ \midrule
1 & Classical description of the atom\\
2 & Quantum model of the atom\\
3 & Quantum numbers and atomic orbitals\\
4 & Electronic configuration and periodic table\\
5 & Atomic properties and periodicity\\ \bottomrule
\end{tabular}
\end{center}
